\section{Chatbot trong hỗ trợ tra cứu thủ tục
hành
chính}

Chatbot là một phần mềm tiếp nhận các yêu cầu của người dùng sau đó phản
hồi theo dạng hội thoại. Trong các hệ thống thông thường, chatbot có thể
dùng để chăm sóc khách hàng, trả lời câu hỏi lặp lại hoặc hướng dẫn thao
tác. Trong hành chính công, yêu cầu đặt ra chặt chẽ hơn vì thông tin trả
lời liên quan trực tiếp đến giấy tờ, thời hạn, phí, lệ phí và các cơ
quan có thẩm quyền.

Người dân thường không nhớ chính xác tên thủ tục hoặc mã thủ tục. Họ mô
tả nhu cầu bằng lời nói hằng ngày, ví dụ ``làm giấy khai sinh cần gì'',
``đăng ký kết hôn ở đâu'', ``cấp lại giấy tờ mất bao lâu''. Vì vậy,
chatbot trong đề tài không chỉ là một ô tìm kiếm có giao diện hội thoại,
mà là lớp trung gian giúp chuyển câu hỏi tự nhiên thành truy vấn có thể
xử lý trên kho dữ liệu thủ tục.

Trong phạm vi đồ án, chatbot được xem là kênh hỗ trợ tra cứu ban đầu. Hệ
thống không thay thế cán bộ chuyên môn và không đưa ra quyết định hành
chính. Vai trò của hệ thống là tìm thông tin phù hợp, diễn giải ngắn gọn
và thông báo rõ khi dữ liệu hiện tại không có căn cứ để trả lời.

Có hai hướng chatbot phổ biến. Chatbot dựa trên luật dễ kiểm soát nhưng
kém linh hoạt khi người dùng thay đổi cách diễn đạt. Chatbot dùng LLM
linh hoạt hơn nhưng có nguy cơ sinh thông tin không có trong dữ liệu
nguồn. Với bài toán TTHC, hướng phù hợp là chatbot AI có ràng buộc bởi
kho dữ liệu đã kiểm soát thông qua kỹ thuật RAG.

\section{Xử lý ngôn ngữ tự nhiên và mô hình
ngôn ngữ
lớn}

NLP giúp hệ thống tiếp nhận và phân tích câu hỏi tiếng Việt của người
dùng. Trong bài toán này, xử lý ngôn ngữ không dừng ở việc tách từ hay
tìm kiếm từ khóa. Hệ thống cần nhận biết các cách nói khác nhau của cùng
một ý định, chẳng hạn ``cần giấy tờ gì'', ``hồ sơ gồm gì'', ``chuẩn bị
những gì'' đều hướng đến nhóm thông tin thành phần hồ sơ.

LLM hỗ trợ đọc và diễn giải các đoạn dữ liệu hành chính dài thành câu
trả lời dễ hiểu hơn. Tuy nhiên, mô hình không được sử dụng như nguồn tri
thức độc lập. Nếu để mô hình tự trả lời theo kiến thức sẵn có, hệ thống
có thể sinh thông tin không đúng với thủ tục đang áp dụng. Vì vậy, mô
hình trong đề tài chỉ đóng vai trò tổng hợp và trình bày lại context đã
được truy xuất từ kho dữ liệu.

Để giảm rủi ro, prompt gửi cho mô hình cần quy định rõ: chỉ trả lời theo
dữ liệu được cung cấp; nếu dữ liệu không đề cập thì phải thông báo không
có thông tin trong kho hiện tại; không tự suy diễn mức phí, thời hạn
hoặc giấy tờ. Nguyên tắc này đặc biệt quan trọng với hệ thống hành chính
công vì người dùng có thể dựa vào câu trả lời để chuẩn bị hồ sơ.

Trong mã nguồn của đồ án, mô hình sinh câu trả lời được cấu hình theo cơ
chế model chính và fallback. Giai đoạn phát triển ưu tiên mô hình chạy
cục bộ qua Ollama để giảm chi phí thử nghiệm, đồng thời có thể chuyển
sang các dịch vụ khác khi mô hình chính không khả dụng. Thiết kế này
giúp hệ thống linh hoạt hơn trong quá trình chạy thử và demo.

\section{Embedding, tìm kiếm ngữ nghĩa và
reranking}

Embedding là bước chuyển văn bản thành vector số. Khi câu hỏi và đoạn
tài liệu đều được biểu diễn dưới dạng vector, hệ thống có thể so sánh độ
gần nhau về mặt ngữ nghĩa thay vì chỉ so khớp chuỗi ký tự. Đây là nền
tảng để chatbot hiểu được các câu hỏi có cách diễn đạt khác nhau nhưng
cùng ý nghĩa.

Trong đồ án, mô hình embedding được sử dụng là
sentence-transformers/all-MiniLM-L6-v2 thông qua lớp
HuggingFaceEmbeddings. Mô hình này chạy ở phía backend theo hướng local:
khi khởi tạo, hệ thống tải mô hình về môi trường chạy và dùng trực tiếp
để tạo vector cho dữ liệu cũng như câu hỏi. Như vậy, bước embedding
không phụ thuộc vào API bên ngoài cho từng truy vấn, phù hợp với giai
đoạn prototype và thử nghiệm.

Ưu điểm của all-MiniLM-L6-v2 là nhẹ, tốc độ xử lý nhanh và đủ phù hợp
cho bài toán truy xuất văn bản ngắn đến trung bình. Tuy nhiên, mô hình
này không phải mô hình tối ưu riêng cho tiếng Việt pháp lý. Vì vậy, hệ
thống bổ sung bước chuẩn hóa synonym, nhận diện tên thủ tục và reranking
để giảm sai lệch trong quá trình chọn context.

Sau khi semantic search trả về danh sách chunk ứng viên, hệ thống sử
dụng CrossEncoder với mô hình cross-encoder/ms-marco-MiniLM-L-6-v2 để
chấm lại mức độ phù hợp giữa câu hỏi và từng chunk. Khác với embedding
search chỉ so sánh vector đã tạo sẵn, CrossEncoder đánh giá trực tiếp
cặp câu hỏi - tài liệu, do đó có khả năng xếp lại các kết quả gần nghĩa
nhưng chưa đúng trọng tâm. Nếu reranker không tải được, hệ thống vẫn có
thể chạy theo kết quả truy xuất ban đầu, nhưng chất lượng chọn context
có thể giảm.

\section{Kỹ thuật Retrieval-Augmented
Generation
(RAG)}

RAG là một kỹ thuật kết hợp giữa việc truy xuất thông tin và tạo ra văn
bản. Hệ thống trước tiên tìm các đoạn dữ liệu liên quan từ kho tri thức,
sau đó đưa các đoạn này vào ngữ cảnh để mô hình ngôn ngữ sinh câu trả
lời. Cách làm này phù hợp với những lĩnh vực cần dữ liệu có nguồn rõ
ràng và thường xuyên cập nhật.

Trong bài toán TTHC, RAG tách hai nhiệm vụ thành hai pha. Pha truy xuất
chịu trách nhiệm tìm đúng thủ tục và đúng trường thông tin như hồ sơ,
phí, lệ phí, thời hạn hoặc cơ quan thực hiện. Pha sinh chịu trách nhiệm
diễn giải dữ liệu được tìm thấy thành câu trả lời ngắn gọn, dễ đọc. Nhờ
vậy, mô hình ngôn ngữ không phải tự nhớ toàn bộ dữ liệu pháp lý.

Pha xây dựng kho tri thức bắt đầu từ dữ liệu thủ tục đã chuẩn hóa. Mỗi
thủ tục được chia thành các chunk theo ý nghĩa nghiệp vụ thay vì cắt cơ
học theo độ dài. Ví dụ, thông tin tổng quan, cách thức thực hiện, thành
phần hồ sơ và căn cứ pháp lý được tách thành những nhóm riêng. Với các
đoạn quá dài, hệ thống mới tiếp tục chia nhỏ bằng
RecursiveCharacterTextSplitter để đảm bảo kích thước phù hợp cho truy
xuất.

Pha hỏi đáp bắt đầu khi người dùng nhập câu hỏi. Hệ thống có thể xử lý
intent xã giao, viết lại câu hỏi dựa trên lịch sử, chuẩn hóa từ đồng
nghĩa, truy xuất dữ liệu trong ChromaDB, rerank kết quả và chọn context
cuối cùng. Context này được đưa vào prompt cho LLM sinh câu trả lời. Nếu
context không chứa thông tin cần hỏi, câu trả lời phải thể hiện rằng dữ
liệu hiện tại không đề cập.

\begin{longtable}[]{|
  >{\raggedright\arraybackslash}p{(\columnwidth - 4\tabcolsep) * \real{0.1716}}|
  >{\raggedright\arraybackslash}p{(\columnwidth - 4\tabcolsep) * \real{0.3128}}|
  >{\raggedright\arraybackslash}p{(\columnwidth - 4\tabcolsep) * \real{0.5156}}|}
\caption{Vai trò các thành phần trong kỹ thuật RAG}\\
\hline
\begin{minipage}[b]{\linewidth}\raggedright
\textbf{Thành phần}
\end{minipage} & \begin{minipage}[b]{\linewidth}\raggedright
\textbf{Vai trò}
\end{minipage} & \begin{minipage}[b]{\linewidth}\raggedright
\textbf{Ý nghĩa trong đề tài}
\end{minipage} \\
\hline
\endfirsthead
\continuedcaption{Vai trò các thành phần trong kỹ thuật RAG}
\hline
\begin{minipage}[b]{\linewidth}\raggedright
\textbf{Thành phần}
\end{minipage} & \begin{minipage}[b]{\linewidth}\raggedright
\textbf{Vai trò}
\end{minipage} & \begin{minipage}[b]{\linewidth}\raggedright
\textbf{Ý nghĩa trong đề tài}
\end{minipage} \\
\hline
\endhead
\hline
\endlastfoot
Dữ liệu nguồn & Lưu thông tin thủ tục & Là căn cứ để chatbot trả lời,
hạn chế suy diễn. \\
\hline
Chunking & Chia dữ liệu theo nhóm ý nghĩa & Giúp truy xuất đúng phần hồ
sơ, phí, thời hạn hoặc cơ quan. \\
\hline
Embedding & Biểu diễn văn bản thành vector & Cho phép tìm kiếm theo ngữ
nghĩa, không phụ thuộc hoàn toàn vào từ khóa. \\
\hline
Vector database & Lưu vector và metadata & Hỗ trợ tìm nhanh các chunk
liên quan đến câu hỏi. \\
\hline
Reranker & Xếp hạng lại kết quả truy xuất & Giảm nguy cơ chọn nhầm thủ
tục hoặc nhầm trường thông tin. \\
\hline
LLM & Sinh câu trả lời từ context & Diễn giải dữ liệu hành chính thành
câu trả lời dễ hiểu. \\
\hline
\end{longtable}

\section{Vector database và
ChromaDB}

Vector database là cơ sở dữ liệu được thiết kế để lưu trữ vector
embedding và tìm kiếm các vector gần nhau theo một thước đo tương đồng.
Trong hệ thống RAG, vector database đóng vai trò cầu nối giữa câu hỏi
của người dùng và kho dữ liệu thủ tục đã được chia chunk.

Đồ án sử dụng ChromaDB để lưu các vector embedding kèm metadata.
Metadata có thể gồm mã thủ tục, tên thủ tục, lĩnh vực, loại chunk và các
thông tin nhận diện khác. Nhờ metadata, hệ thống không chỉ tìm theo độ
tương đồng ngữ nghĩa mà còn có thể lọc theo thủ tục hoặc nhóm thông tin
khi cần. Điều này hữu ích vì nhiều TTHC có tên gần giống nhau hoặc cùng
thuộc một lĩnh vực.

ChromaDB phù hợp với phạm vi prototype và đồ án vì dễ cài đặt, dễ chạy
cục bộ, có khả năng persist dữ liệu xuống thư mục và tích hợp thuận tiện
với LangChain. Trong hệ thống hiện tại, thư mục chroma\_db được sử dụng
để lưu kho vector; khi dữ liệu thủ tục thay đổi, quản trị viên có thể
build lại ChromaDB từ file JSON đã cập nhật.

Tuy nhiên, nếu triển khai ở quy mô lớn, ChromaDB cục bộ cần được đánh
giá lại về khả năng mở rộng, đồng thời và vận hành. Khi số lượng thủ
tục, số lượng người dùng hoặc yêu cầu sẵn sàng tăng cao, cần xem xét các
giải pháp như vector database triển khai dạng server, dịch vụ managed
vector search, cơ chế replication, backup tự động, giám sát tài nguyên
và phân tách môi trường cập nhật dữ liệu với môi trường phục vụ truy
vấn.

\section{Đánh giá chất lượng hệ thống
RAG}

Đánh giá hệ thống RAG cần xem xét cả khâu truy xuất và khâu sinh câu trả
lời. Nếu truy xuất sai tài liệu, mô hình ngôn ngữ dù diễn đạt tốt vẫn có
thể trả lời sai. Ngược lại, nếu truy xuất đúng nhưng prompt không chặt
chẽ, mô hình vẫn có thể diễn giải thiếu hoặc thêm thông tin ngoài
context. Vì vậy, việc đánh giá không nên chỉ dừng ở cảm nhận câu trả lời
có tự nhiên hay không.

Trong phạm vi đồ án, có thể xây dựng bộ câu hỏi thử nghiệm có đáp án
tham chiếu từ dữ liệu nguồn. Mỗi câu hỏi được gắn nhóm thông tin cần
kiểm tra như hồ sơ, thời hạn, phí, cơ quan thực hiện, căn cứ pháp lý
hoặc câu hỏi tiếp nối. Khi chạy thử, cần ghi nhận chunk được truy xuất,
câu trả lời cuối cùng và thời gian phản hồi để phân tích lỗi.

\begin{longtable}[]{|
  >{\raggedright\arraybackslash}p{(\columnwidth - 4\tabcolsep) * \real{0.2054}}|
  >{\raggedright\arraybackslash}p{(\columnwidth - 4\tabcolsep) * \real{0.3728}}|
  >{\raggedright\arraybackslash}p{(\columnwidth - 4\tabcolsep) * \real{0.4219}}|}
\caption{Một số tiêu chí đánh giá hệ thống RAG}\\
\hline
\begin{minipage}[b]{\linewidth}\raggedright
\textbf{Tiêu chí}
\end{minipage} & \begin{minipage}[b]{\linewidth}\raggedright
\textbf{Ý nghĩa}
\end{minipage} & \begin{minipage}[b]{\linewidth}\raggedright
\textbf{Cách áp dụng trong đồ án}
\end{minipage} \\
\hline
\endfirsthead
\continuedcaption{Một số tiêu chí đánh giá hệ thống RAG}
\hline
\begin{minipage}[b]{\linewidth}\raggedright
\textbf{Tiêu chí}
\end{minipage} & \begin{minipage}[b]{\linewidth}\raggedright
\textbf{Ý nghĩa}
\end{minipage} & \begin{minipage}[b]{\linewidth}\raggedright
\textbf{Cách áp dụng trong đồ án}
\end{minipage} \\
\hline
\endhead
\hline
\endlastfoot
Độ chính xác truy xuất Top-k & Đánh giá các chunk liên quan có nằm trong
nhóm kết quả đầu hay không & So sánh context truy xuất với đáp án chuẩn
của từng câu hỏi. \\
\hline
Độ đúng của câu trả lời & Đánh giá câu trả lời có khớp dữ liệu nguồn
không & Kiểm tra thủ công theo trường hồ sơ, thời hạn, phí, cơ quan. \\
\hline
Mức độ bám nguồn & Đánh giá mô hình có thêm thông tin ngoài context
không & Đánh dấu các câu trả lời suy diễn hoặc không có trong dữ
liệu. \\
\hline
Tỉ lệ không trả lời khi thiếu dữ liệu & Đánh giá khả năng từ chối đúng
khi kho dữ liệu không có thông tin & Dùng câu hỏi ngoài phạm vi và kiểm
tra phản hồi ``dữ liệu hiện tại không đề cập''. \\
\hline
Latency & Đánh giá thời gian phản hồi toàn tuyến & Đo thời gian từ lúc
gửi câu hỏi đến lúc nhận câu trả lời. \\
\hline
Độ ổn định & Đánh giá khả năng fallback khi mô hình chính lỗi & Kiểm tra
các trường hợp timeout hoặc model không phản hồi. \\
\hline
\end{longtable}

Các tiêu chí trên giúp việc đánh giá có cơ sở hơn so với nhận xét chung
chung. Trong bối cảnh hành chính công, hai tiêu chí cần ưu tiên là độ
đúng theo dữ liệu nguồn và mức độ bám nguồn. Một câu trả lời ngắn nhưng
đúng và có kiểm soát tốt hơn một câu trả lời dài, mạch lạc nhưng thêm
thông tin không có trong thủ tục.

\section{Công nghệ web và lưu trữ được áp
dụng với hệ
thống}

Về phần mềm ứng dụng, hệ thống sử dụng kiến trúc tách biệt frontend và
backend. Backend được xây dựng bằng Python và FastAPI để điều phối API,
xác thực, truy xuất ChromaDB, gọi pipeline RAG và lưu lịch sử. Python
phù hợp với đề tài vì có hệ sinh thái thư viện AI và xử lý dữ liệu phong
phú.

Frontend được xây dựng bằng React, TypeScript, Vite và Tailwind CSS.
React giúp chia giao diện thành các component như ChatBox, Sidebar,
Message, InputBox, LoginPage và AdminDashboard. TypeScript hỗ trợ kiểm
soát kiểu dữ liệu, Vite giúp chạy thử nhanh, còn Tailwind CSS giúp xây
dựng giao diện responsive cho cả máy tính và thiết bị di động.

Firebase Authentication được dùng cho đăng nhập và cấp ID Token. Backend
dùng Firebase Admin SDK để xác minh token trước khi xử lý các API cần
bảo vệ. Firestore lưu dữ liệu vận hành như phiên chat, tin nhắn, danh
sách admin, mã OTP và cấu hình backup. Như vậy, Firebase đảm nhận dữ
liệu nghiệp vụ, còn ChromaDB đảm nhận dữ liệu tri thức phục vụ RAG.

Tóm lại, phần công nghệ web trong đồ án giữ vai trò triển khai sản phẩm,
còn phần cốt lõi quyết định chất lượng trả lời nằm ở dữ liệu, embedding,
truy xuất, reranking, prompt và cách đánh giá RAG. Vì vậy, báo cáo tập
trung nhiều hơn vào quy trình RAG và yêu cầu hành chính công thay vì mô
tả quá dài các công nghệ giao diện.
